% TOP_PROBLEM: {"id": 30007086, "title": "Fractional chromatic number versus ordinary Hadwiger number", "category_id": 2, "category": {"id": 2, "name": "combinatorics", "display_name": "Combinatorics", "description": "Counting problems, graph theory, discrete structures.", "slug": "combinatorics", "order_index": 2, "created_at": "2026-07-31T15:26:25.671Z"}, "status": "open"}
% FIRST_POSTED: 2026-10-06
% ATTEMPT_STATUS: solved
% OPENAI_CLAIM: {"schema_version":1,"source_repo":"https://github.com/openai/math","source_commit":"adc7f1241b42e322a6451854ab7e4b4c146bf78a","release_date":"2026-10-06","match":"full","resolution":"disproved","independently_reviewed":false,"second_pass":"agreed","families":[{"family":"157","title":"Graph coloring, clique minors, and Colin de Verdière invariants","subject":"Combinatorics","manuscripts":[{"title":"A counterexample to Hadwiger's conjecture","pdf_path":"preprints/A-counterexample-to-Hadwigers-conjecture-September-23-2026/paper.pdf","readme_path":"preprints/A-counterexample-to-Hadwigers-conjecture-September-23-2026/README.md","date":"2026-09-23","theorem_ref":"thm:main (sections/01-introduction.tex:20); cor:hadwiger (sections/01-introduction.tex:48)","statement_sources":[{"path":"preprints/A-counterexample-to-Hadwigers-conjecture-September-23-2026/build/sections/01-introduction.tex","line":20,"label":"thm:main"},{"path":"preprints/A-counterexample-to-Hadwigers-conjecture-September-23-2026/build/sections/01-introduction.tex","line":48,"label":"cor:hadwiger"}]},{"title":"A counterexample to the Colin de Verdière chromatic conjecture","pdf_path":"preprints/A-counterexample-to-the-Colin-de-Verdiere-chromatic-conjecture-September-23-2026/paper.pdf","readme_path":"preprints/A-counterexample-to-the-Colin-de-Verdiere-chromatic-conjecture-September-23-2026/README.md","date":"2026-09-23","theorem_ref":"thm:main (sections/01-introduction.tex:42); thm:base-graphs (sections/01-introduction.tex:147); cor:fractional-spectral (sections/01-introduction.tex:221); cor:list-spectral (sections/01-introduction.tex:262)","statement_sources":[{"path":"preprints/A-counterexample-to-the-Colin-de-Verdiere-chromatic-conjecture-September-23-2026/build/sections/01-introduction.tex","line":42,"label":"thm:main"},{"path":"preprints/A-counterexample-to-the-Colin-de-Verdiere-chromatic-conjecture-September-23-2026/build/sections/01-introduction.tex","line":147,"label":"thm:base-graphs"},{"path":"preprints/A-counterexample-to-the-Colin-de-Verdiere-chromatic-conjecture-September-23-2026/build/sections/01-introduction.tex","line":221,"label":"cor:fractional-spectral"},{"path":"preprints/A-counterexample-to-the-Colin-de-Verdiere-chromatic-conjecture-September-23-2026/build/sections/01-introduction.tex","line":262,"label":"cor:list-spectral"}]}],"lean":{"doc_path":"lean/docs/157.md","comparators":[],"declarations":[],"covers_main_theorem":false},"reasoning_trace":null}]}
\documentclass[11pt]{article}
\usepackage[margin=25mm]{geometry}
\usepackage{fontspec}
\setmainfont{Latin Modern Roman}
\usepackage{amsmath,amssymb,mathtools}
\usepackage{unicode-math}
\setmathfont{Latin Modern Math}
\usepackage{enumitem}
\usepackage{xurl,hyperref,bookmark}
\hypersetup{colorlinks=true,urlcolor=blue}
\setcounter{secnumdepth}{0}
\setlength{\parindent}{0pt}
\setlength{\parskip}{5pt}
\setlength{\emergencystretch}{3em}
\begin{document}
\section*{Fractional chromatic number versus ordinary Hadwiger number}

\subsection{OpenAI's claim}
\textit{Chromatic number beyond clique-minor size}: The counterexample paper claims arbitrarily large finite graphs with $\chi_f(G)>h(G)$, hence also $\chi(G)>h(G)$; no independent vertex set in an example contains three vertices.

\textit{Colin de Verdière chromatic obstruction}: The manuscript claims finite graphs whose fractional chromatic number is greater than $\mu(G)+1$, and therefore whose ordinary chromatic number violates the same bound; such examples exist at arbitrarily large orders.

\subsection{Mathematical statement of the claimed result}
\paragraph{Mathematical claim 1.} There are finite simple graphs $G$ of arbitrarily large order $m$ with independence number $\alpha(G)\le2$ and maximum size of a pairwise-touching matching less than $m/100$. The resulting corollary states \[h(G)<26m/75+2/3<m/2\le\chi_f(G)\le\chi(G),\] where $h(G)=\max\{t:K_t\text{ is a minor of }G\}$ and $\chi_f,\chi$ are fractional and ordinary chromatic numbers. Thus the exact inequalities $\chi\le h$ and $\chi_f\le h$ fail; no conclusion about replacing $h$ by a fractional minor parameter is supplied.

\paragraph{Mathematical claim 2.} The claimed counterexamples are finite simple graphs of arbitrarily large order $m$ with \[\alpha(G)\le2,\qquad\mu(G)+1<m/2\le\chi(G),\] where $\mu(G)$ is the Colin de Verdi\`ere spectral graph invariant. The fractional-coloring corollary strengthens the obstruction to \[\chi_f(G)\ge m/\alpha(G)\ge m/2>\kappa(G)+1\ge\mu(G)+1\ge h(G).\] Here $\kappa$ is its strong-Arnold-property-free spectral relaxation, and $h$ is the largest complete-minor order. In particular both $\chi\le\mu+1$ and $\chi_f\le h$ are refuted. A further corollary gives $\chi_{\mathrm{list}}(G)\le C(\mu(G)+1)$ with an absolute $C$.

\subsection{Outline of the argument}
\textit{Chromatic number beyond clique-minor size}: The graph construction uses Boolean moment representations and paired frames to control compatible vertex distributions. Entropy projection, fingerprints for independent sets, and a capacitated-cut argument transfer the distributional constraints to finite graphs. The resulting obstruction limits clique-minor size, while the small independence number gives a larger fractional chromatic number.

\textit{Colin de Verdière chromatic obstruction}: Allowed channel images are thinned while preserving the primal distribution, supplying enough compatible pairs for the conflict estimates. A Boolean moment argument restricts low-rank matrix representations. Entropy and finite sampling then produce graphs for which every well-signed matrix with a single negative eigenvalue has too much rank to support the proposed chromatic inequality.

\subsection{Scope relative to this problem}
\textit{Chromatic number beyond clique-minor size}: The named fractional-coloring corollary displays chi\_f(G)>h(G) for arbitrarily large finite nonempty simple graphs. It refutes the relevant assertion using the ordinary complete-minor number h, and does not require a fractional-minor definition.

\textit{Colin de Verdière chromatic obstruction}: The named fractional-coloring corollary displays chi\_f(G)>h(G) for arbitrarily large finite nonempty simple graphs. It refutes the relevant assertion using the ordinary complete-minor number h, and does not require a fractional-minor definition.

The claimed counterexample lies within the question’s hypotheses and contradicts its universal conclusion.

\subsection{Formal verification}
Family 157 describes Lean coverage for the list-coloring upper bound with a universal multiple of clique-minor size; neither chromatic counterexample is part of this selection.

\textit{Chromatic number beyond clique-minor size}: No manuscript-specific Comparator is listed in the pinned inventory.

The documented selection does not cover the entire manuscript assertion relevant here.

\textit{Colin de Verdière chromatic obstruction}: No manuscript-specific Comparator is listed in the pinned inventory.

The documented selection does not cover the entire manuscript assertion relevant here.

This site has not run the Lean development. These links describe OpenAI’s selected formal statements and should be read with their stated scope.

\subsection{Status caveat}
OpenAI reports that an internal model produced the manuscripts. Their arguments have not been independently verified as proofs by this site. Any portion outside the documented formal scope may contain mistakes, and the catalogue label reports the announced mathematical claim.

\subsection{Sources}
\begin{itemize}
\item OpenAI, \emph{A counterexample to Hadwiger's conjecture}, 2026-09-23; citation key \texttt{OAI:A-counterexample-to-Hadwigers-conjecture-September-23-2026}. \href{https://github.com/openai/math/blob/adc7f1241b42e322a6451854ab7e4b4c146bf78a/preprints/A-counterexample-to-Hadwigers-conjecture-September-23-2026/paper.pdf}{PDF (GitHub)}; \href{https://github.com/openai/math/raw/adc7f1241b42e322a6451854ab7e4b4c146bf78a/preprints/A-counterexample-to-Hadwigers-conjecture-September-23-2026/paper.pdf}{PDF (direct)}; \href{https://github.com/openai/math/blob/adc7f1241b42e322a6451854ab7e4b4c146bf78a/preprints/A-counterexample-to-Hadwigers-conjecture-September-23-2026/README.md}{Manuscript page and citation}; \href{https://github.com/openai/math/blob/main/preprints/A-counterexample-to-Hadwigers-conjecture-September-23-2026/paper.pdf}{Latest version}; Statement source: \href{https://github.com/openai/math/blob/adc7f1241b42e322a6451854ab7e4b4c146bf78a/preprints/A-counterexample-to-Hadwigers-conjecture-September-23-2026/build/sections/01-introduction.tex\#L20}{thm:main}; Statement source: \href{https://github.com/openai/math/blob/adc7f1241b42e322a6451854ab7e4b4c146bf78a/preprints/A-counterexample-to-Hadwigers-conjecture-September-23-2026/build/sections/01-introduction.tex\#L48}{cor:hadwiger}.
\item OpenAI, \emph{A counterexample to the Colin de Verdière chromatic conjecture}, 2026-09-23; citation key \texttt{OAI:A-counterexample-to-the-Colin-de-Verdiere-chromatic-conjecture-September-23-2026}. \href{https://github.com/openai/math/blob/adc7f1241b42e322a6451854ab7e4b4c146bf78a/preprints/A-counterexample-to-the-Colin-de-Verdiere-chromatic-conjecture-September-23-2026/paper.pdf}{PDF (GitHub)}; \href{https://github.com/openai/math/raw/adc7f1241b42e322a6451854ab7e4b4c146bf78a/preprints/A-counterexample-to-the-Colin-de-Verdiere-chromatic-conjecture-September-23-2026/paper.pdf}{PDF (direct)}; \href{https://github.com/openai/math/blob/adc7f1241b42e322a6451854ab7e4b4c146bf78a/preprints/A-counterexample-to-the-Colin-de-Verdiere-chromatic-conjecture-September-23-2026/README.md}{Manuscript page and citation}; \href{https://github.com/openai/math/blob/main/preprints/A-counterexample-to-the-Colin-de-Verdiere-chromatic-conjecture-September-23-2026/paper.pdf}{Latest version}; Statement source: \href{https://github.com/openai/math/blob/adc7f1241b42e322a6451854ab7e4b4c146bf78a/preprints/A-counterexample-to-the-Colin-de-Verdiere-chromatic-conjecture-September-23-2026/build/sections/01-introduction.tex\#L42}{thm:main}; Statement source: \href{https://github.com/openai/math/blob/adc7f1241b42e322a6451854ab7e4b4c146bf78a/preprints/A-counterexample-to-the-Colin-de-Verdiere-chromatic-conjecture-September-23-2026/build/sections/01-introduction.tex\#L147}{thm:base-graphs}; Statement source: \href{https://github.com/openai/math/blob/adc7f1241b42e322a6451854ab7e4b4c146bf78a/preprints/A-counterexample-to-the-Colin-de-Verdiere-chromatic-conjecture-September-23-2026/build/sections/01-introduction.tex\#L221}{cor:fractional-spectral}; Statement source: \href{https://github.com/openai/math/blob/adc7f1241b42e322a6451854ab7e4b4c146bf78a/preprints/A-counterexample-to-the-Colin-de-Verdiere-chromatic-conjecture-September-23-2026/build/sections/01-introduction.tex\#L262}{cor:list-spectral}.
\item Family 157: \href{https://github.com/openai/math/blob/adc7f1241b42e322a6451854ab7e4b4c146bf78a/lean/docs/157.md}{Lean scope}.
\item \href{https://github.com/openai/math/blob/adc7f1241b42e322a6451854ab7e4b4c146bf78a/CONTENTS.md}{OpenAI catalogue entry}.
\end{itemize}
\end{document}
